% ═══════════════════════════════════════════════════════════════
% LEHRERHINWEISE: El tiempo y los lugares
% Spanisch Klasse 8 - Unidad 6, DS 1
% ═══════════════════════════════════════════════════════════════

\documentclass[11pt, a4paper]{article}

\usepackage[utf8]{inputenc}
\usepackage[T1]{fontenc}
\usepackage[ngerman]{babel}
\usepackage{helvet}
\renewcommand{\familydefault}{\sfdefault}

\usepackage[margin=2cm]{geometry}
\usepackage{enumitem}
\usepackage{tabularx}
\usepackage{booktabs}
\usepackage{longtable}

\usepackage{xcolor}
\usepackage{amssymb}
\usepackage{tcolorbox}
\tcbuselibrary{skins}

\usepackage{fancyhdr}
\usepackage{lastpage}
\usepackage{needspace}

% ═══════════════════════════════════════════════════════════════
% FARBEN
% ═══════════════════════════════════════════════════════════════

\definecolor{spanisch-color}{HTML}{c41e3a}
\definecolor{grau-color}{HTML}{888888}
\definecolor{hellgrau-color}{HTML}{f5f5f5}
\definecolor{basis-color}{HTML}{2a7a4b}
\definecolor{standard-color}{HTML}{2c5aa0}
\definecolor{erweiterung-color}{HTML}{7b1fa2}
\definecolor{loesung-color}{HTML}{c62828}

\colorlet{spanisch}{spanisch-color}
\colorlet{grau}{grau-color}
\colorlet{hellgrau}{hellgrau-color}
\colorlet{basis}{basis-color}
\colorlet{standard}{standard-color}
\colorlet{erweiterung}{erweiterung-color}
\colorlet{loesung}{loesung-color}

\newcommand{\fachfarbe}{spanisch}

% ═══════════════════════════════════════════════════════════════
% VARIABLEN
% ═══════════════════════════════════════════════════════════════

\newcommand{\thema}{El tiempo y los lugares}
\newcommand{\sequenz}{06}
\newcommand{\block}{a}
\newcommand{\fach}{Spanisch}
\newcommand{\klasse}{8}
\newcommand{\dauer}{90 Min. (Doppelstunde)}
\newcommand{\lerngruppengroesse}{24}
\newcommand{\datum}{\today}

% ═══════════════════════════════════════════════════════════════
% KOPF- UND FUSSZEILE
% ═══════════════════════════════════════════════════════════════

\pagestyle{fancy}
\fancyhf{}
\renewcommand{\headrulewidth}{0.4pt}
\renewcommand{\footrulewidth}{0.4pt}

\lhead{\fach{} -- Klasse \klasse}
\chead{\textbf{LH-\sequenz\block{} (Lehrerhinweise)}}
\rhead{\datum}

\lfoot{}
\cfoot{}
\rfoot{Seite \thepage\ von \pageref{LastPage}}

% ═══════════════════════════════════════════════════════════════
% BEFEHLE
% ═══════════════════════════════════════════════════════════════

\newcommand{\B}{\textcolor{basis}{\textbf{[B]}}}
\newcommand{\Std}{\textcolor{standard}{\textbf{[S]}}}
\newcommand{\E}{\textcolor{erweiterung}{\textbf{[E]}}}
\newcommand{\lsg}[1]{\textcolor{loesung}{\textbf{#1}}}
\newcommand{\es}[1]{\textcolor{\fachfarbe}{\textbf{#1}}}

\newtcolorbox{kurzplan}{
    colback=hellgrau,
    colframe=\fachfarbe,
    colbacktitle=white,
    coltitle=\fachfarbe,
    boxrule=1pt,
    arc=0pt,
    fonttitle=\bfseries,
    attach boxed title to top left={yshift=-2mm, xshift=5mm},
    boxed title style={boxrule=0pt, colback=white},
    title=Kurzplan
}

\newtcolorbox{differenzierung}{
    colback=white,
    colframe=grau,
    colbacktitle=white,
    coltitle=grau,
    boxrule=0.5pt,
    arc=0pt,
    fonttitle=\bfseries,
    attach boxed title to top left={yshift=-2mm, xshift=5mm},
    boxed title style={boxrule=0pt, colback=white},
    title=Differenzierung
}

\newtcolorbox{fehlerbox}{
    colback=white,
    colframe=loesung,
    colbacktitle=white,
    coltitle=loesung,
    boxrule=0.5pt,
    arc=0pt,
    fonttitle=\bfseries,
    attach boxed title to top left={yshift=-2mm, xshift=5mm},
    boxed title style={boxrule=0pt, colback=white},
    title=Häufige Fehler
}

\newtcolorbox{materialbox}{
    colback=white,
    colframe=\fachfarbe,
    colbacktitle=white,
    coltitle=\fachfarbe,
    boxrule=0.5pt,
    arc=0pt,
    fonttitle=\bfseries,
    attach boxed title to top left={yshift=-2mm, xshift=5mm},
    boxed title style={boxrule=0pt, colback=white},
    title=Material-Checkliste
}

\newtcolorbox{spielebox}{
    colback=white,
    colframe=basis,
    colbacktitle=white,
    coltitle=basis,
    boxrule=0.5pt,
    arc=0pt,
    fonttitle=\bfseries,
    attach boxed title to top left={yshift=-2mm, xshift=5mm},
    boxed title style={boxrule=0pt, colback=white},
    title=Spielvorschläge (ad-hoc)
}

% ═══════════════════════════════════════════════════════════════
% DOKUMENT
% ═══════════════════════════════════════════════════════════════

\begin{document}

% ═══════════════════════════════════════════════════════════════
% KOPFBEREICH
% ═══════════════════════════════════════════════════════════════

\begin{center}
    {\Large\bfseries\textcolor{\fachfarbe}{Lehrerhinweise: \thema}}\\[0.3em]
    {\normalsize LH-\sequenz\block{} | \dauer{} | Qué pasa 1, S. 102-103}
\end{center}

\vspace{10pt}

% ═══════════════════════════════════════════════════════════════
% 1. KURZPLAN
% ═══════════════════════════════════════════════════════════════

\section*{1. Stundenverlauf (90 Min.)}

\begin{tabularx}{\textwidth}{|c|l|X|l|}
\hline
\textbf{Zeit} & \textbf{Phase} & \textbf{Inhalt} & \textbf{Material} \\
\hline
0--10 & Einstieg & L zeigt 6 Fotos (Buch S. 102-103), fragt: ¿Qué veis? SuS beschreiben frei auf Deutsch/Spanisch & Buch \\
\hline
10--25 & Input Wetter & L führt Wetterausdrücke ein mit Gestik/Bildern. Chorisches Wiederholen. $\rightarrow$ \textbf{LB Kasten 1} ausfüllen & LB \\
\hline
25--35 & Übung 1 & Foto-Beschreibung: ¿Qué tiempo hace en la foto 1/2/3...? SuS antworten mündlich & Buch \\
\hline
35--50 & Übung 2 & \textbf{Partner-Interview $\rightarrow$ AB Aufgabe 1:} ¿Qué haces cuando hace sol? & AB \\
\hline
50--60 & Input Orte & L führt Zusatzvokabular (lugares) mit Bildern ein. $\rightarrow$ \textbf{LB Kasten 2} ausfüllen & LB, Bilder \\
\hline
60--75 & Übung 3 & \textbf{Wetter + Kleidung $\rightarrow$ AB Aufgabe 2:} ¿Qué ropa llevas cuando...? Partner-Kette & AB \\
\hline
75--85 & Spiel & \textbf{Wetter-Bingo} oder \textbf{Kettenspiel} (siehe Spielvorschläge) & Bingo-Karten \\
\hline
85--90 & Sicherung & Zusammenfassung: 3 wichtigste Wetterausdrücke + 3 Orte nennen & -- \\
\hline
\end{tabularx}

% ═══════════════════════════════════════════════════════════════
% 2. DIFFERENZIERUNG
% ═══════════════════════════════════════════════════════════════

\section*{2. Differenzierung}

\begin{differenzierung}
\textbf{Legende:} \B{} Basis (BR) | \Std{} Standard (MR) | \E{} Erweiterung (GY)

\vspace{0.5em}

\textbf{WICHTIG:} Diese Marker sind NUR für die Lehrkraft. Auf Schülermaterial erscheinen KEINE Niveaumarkierungen!
\end{differenzierung}

\vspace{1em}

\begin{tabularx}{\textwidth}{|l|X|X|X|}
\hline
& \B{} \textbf{Basis} & \Std{} \textbf{Standard} & \E{} \textbf{Erweiterung} \\
\hline
\textbf{AB Aufg. 1} & Satzmuster vorgeben: "Cuando hace \_\_\_, voy al/a la \_\_\_." & Freie Satzbildung mit Redemittel-Karte & + 2 zusätzliche Wettersituationen \\
\hline
\textbf{AB Aufg. 2} & Nur 2 Sätze (a+b), Kleidung vorgegeben & Alle 4 Sätze, freie Kleidungswahl & + Eigenes Lieblingswetter beschreiben \\
\hline
\textbf{AB Aufg. 3} & 3 Fotos beschreiben & Alle 6 Fotos & + Vergleich: Foto 1 vs. Foto 4 \\
\hline
\end{tabularx}

% ═══════════════════════════════════════════════════════════════
% 3. HÄUFIGE FEHLER
% ═══════════════════════════════════════════════════════════════

\section*{3. Häufige Fehler}

\begin{fehlerbox}
\begin{tabularx}{\textwidth}{|l|X|X|}
\hline
\textbf{Fehler} & \textbf{Beispiel} & \textbf{Reaktion} \\
\hline
*Es hace sol & *"Es hace calor" statt \es{Hace calor} & Kein "es"! Hace ist unpersönlich \\
\hline
*Hace lluvia & *"Hace lluvia" statt \es{Llueve} & Llueve = eigenständiges Verb \\
\hline
Artikelfehler & *"voy a el parque" statt \es{voy al parque} & Erinnerung: a + el = al \\
\hline
Genus lugares & *"la parque", *"el biblioteca" & Artikelliste visualisieren \\
\hline
\end{tabularx}
\end{fehlerbox}

% ═══════════════════════════════════════════════════════════════
% 4. MATERIAL-CHECKLISTE
% ═══════════════════════════════════════════════════════════════

\section*{4. Material-Checkliste}

\begin{materialbox}
\begin{itemize}[label=$\square$]
    \item LB-06a-tiempo-lugares.pdf (\lerngruppengroesse{} Kopien)
    \item AB-06a-tiempo-lugares.pdf (\lerngruppengroesse{} Kopien)
    \item ML-06a-tiempo-lugares.pdf (1× für Lehrkraft)
    \item Lehrwerk: Qué pasa 1, S. 102--103
    \item Bilder: Orte im Viertel (12 Flashcards oder Beamer)
    \item Optional: Bingo-Vorlage (leeres 3×3 Raster)
\end{itemize}
\end{materialbox}

% ═══════════════════════════════════════════════════════════════
% 5. SPIELVORSCHLÄGE
% ═══════════════════════════════════════════════════════════════

\section*{5. Spielvorschläge (ad-hoc, Partnerarbeit)}

\begin{spielebox}
\textbf{1. Wetter-Bingo (10 Min)}
\begin{itemize}
    \item SuS erstellen 3×3 Raster mit 9 Wetterausdrücken (aus Lernblatt)
    \item L/SuS rufen Wettersituationen auf Deutsch $\rightarrow$ SuS kreuzen spanisches Wort an
    \item Wer zuerst eine Reihe hat, ruft "¡Bingo!"
\end{itemize}

\vspace{0.5em}

\textbf{2. Kettenspiel "Cuando..." (10 Min)}
\begin{itemize}
    \item S1: "Cuando hace sol, voy al parque."
    \item S2 muss mit "parque" weiterarbeiten: "En el parque llevo una camiseta."
    \item S3: "La camiseta es roja. Cuando hace frío, llevo..."
    \item Wer stockt, scheidet aus (oder bekommt Hilfskarte)
\end{itemize}

\vspace{0.5em}

\textbf{3. Pantomime-Raten (5-8 Min)}
\begin{itemize}
    \item S1 spielt Wetter pantomimisch vor (friert, schwitzt, Regenschirm...)
    \item Partner rät: "¡Hace frío!" / "¡Llueve!"
    \item Rollentausch nach 3 Wörtern
\end{itemize}
\end{spielebox}

% ═══════════════════════════════════════════════════════════════
% 6. LÖSUNGEN
% ═══════════════════════════════════════════════════════════════

\section*{6. Lösungen AB-\sequenz\block}

\textbf{Kasten 1 + 2:} Siehe LB (Vokabeln zum Übertragen)

\vspace{0.5em}

\textbf{Aufgabe 1: Partner-Interview} (12 Punkte, je 3P pro Antwort)

Mögliche Antworten:
\begin{itemize}
    \item Cuando hace sol: \lsg{Cuando hace sol, voy al parque / a la playa / juego al fútbol.}
    \item Cuando llueve: \lsg{Cuando llueve, voy al cine / me quedo en casa / leo un libro.}
    \item Cuando hace frío: \lsg{Cuando hace frío, me quedo en casa / voy al centro comercial.}
    \item Cuando hace calor: \lsg{Cuando hace calor, voy a la piscina / tomo un helado.}
\end{itemize}

\vspace{0.5em}

\textbf{Aufgabe 2: Kleidung} (12 Punkte, je 3P pro Satz)
\begin{itemize}
    \item a) \lsg{Cuando hace mucho frío, llevo un abrigo, una bufanda, los guantes y el gorro.}
    \item b) \lsg{Cuando hace mucho calor, llevo pantalones cortos, una camiseta y las gafas de sol.}
    \item c) \lsg{Cuando llueve, llevo las botas y el paraguas.}
    \item d) \lsg{Cuando hace buen tiempo, llevo una camiseta, vaqueros y zapatillas.}
\end{itemize}

\vspace{0.5em}

\textbf{Aufgabe 3: Fotos} (4 Punkte)

Flexibel nach Buchfotos. Beispiele: \lsg{Hace sol, Llueve, Nieva, Hace calor, Hay nubes, Hace frío.}

% ═══════════════════════════════════════════════════════════════
% 7. AUSBLICK
% ═══════════════════════════════════════════════════════════════

\section*{7. Ausblick: Folgestunde (Dienstag, 45 Min)}

Die Folgestunde baut auf den \textbf{Ortsvokabeln} dieser Stunde auf:

\begin{itemize}
    \item \textbf{Thema:} Präpositionen des Ortes (enfrente de, al lado de, entre...)
    \item \textbf{Ziel:} Orte räumlich beschreiben ("La panadería está enfrente del parque.")
    \item \textbf{Buch:} S. 104--107 (Übung 2B: Copito, Übung 4: Barrio)
\end{itemize}

\end{document}
